% \begin{table}[t]
% % \footnotesize

% \vspace{-0.2cm}
% \tabcolsep=0.09cm
% \center
% \begin{tabular}{c | c c c c c}
% \toprule
%   Scene Number  & PSNR$\uparrow$ & LPIPS$\downarrow$ & CLIP$\uparrow$ & FID$\downarrow$ & FVD$\downarrow$ \\
%   \hline
%   \hline
%    1 & & & & &  \\
%    4 & & & & & \\
%    8 & & & & & \\ 
%    12 & & & & & \\
% \bottomrule

% \end{tabular}
% \caption{}
% \vspace{-0.2cm}
% \label{table:ablation_data}
% \end{table}

\begin{table}[t]
% \footnotesize

% \vspace{-0.2cm}
\tabcolsep=0.09cm
\center
\begin{tabular}{c | c c | c c }
\toprule
& \multicolumn{2}{c}{GT Comp.} & \multicolumn{2}{c}{HC Comp.}\\
  Strategy  & PSNR$\uparrow$ & LPIPS$\downarrow$ & PSNR$\uparrow$ & LPIPS$\downarrow$   \\
  \hline
  \hline
Random & 17.70 & 0.4847 & 17.07 & 0.3985 \\
  FOV+Random & 19.17 & 0.3825 & 17.47 & 0.3896 \\
  FOV+Non-adj & 20.11 & 0.3075 & \textbf{18.19} & 0.3571 \\
  FOV+Non-adj+Far-space-time & \textbf{20.22} & \textbf{0.3003} & 18.11 & \textbf{0.3414} \\

\bottomrule

\end{tabular}
\caption{\textbf{Ablation of Memory Retrieval Strategy.} The filtering methods of "FOV" and "Non-adj" (where only one frame from continuous frame sequences is selected as a candidate) effectively filter out useless and redundant information, leading to significant improvements in memory capability.}
\vspace{-0.6cm}
\label{table:ablation_retrieval}
\end{table}